% Planung Game-Projekt bis Semesterende (KLW, 2a und 2g)
% Die Klasse wird vor dem \input mit \def\klasse{2a} oder \def\klasse{2g} gesetzt.

\def\klasseA{2a}
\ifx\klasse\klasseA
	\def\zwischenstand{Mi 2.12.}
	\def\abgabe{Mi 6.1.}
	\def\gespraeche{Mi 13.1.\ und Mi 20.1.}
	\def\nachholtermin{Mi 20.1.}
\else
	\def\zwischenstand{Mo 30.11.}
	\def\abgabe{Mo 4.1.}
	\def\gespraeche{Mo 11.1.\ und Mo 18.1.}
	\def\nachholtermin{Mo 18.1.}
\fi

\begin{frame}
	\centering
	\begin{figure}[H]
		\includegraphics[keepaspectratio, width=.5\textwidth, height=\paperheight]{Logos/FreeFlower.pdf}
	\end{figure}
	Informatik \klasse: Programmieren\\
	\textbf{Game-Projekt: Planung bis Semesterende}
\end{frame}

\begin{frame}{Termine}
	\begin{table}
		\begin{tabular}{ll}
			\textbf{Was}                           & \textbf{Wann}  \\\midrule
			Zwischenstand: spielbare Version zeigen & \zwischenstand \\
			Abgabe Game                            & \abgabe        \\
			Einzelgespräche am Code                & \gespraeche    \\
			Nachholtermin für Abwesende            & \nachholtermin{} (zweiter Gesprächstag) \\
		\end{tabular}
	\end{table}
	\vspace{1em}
	Bis zur Abgabe arbeiten Sie in allen Halbklassenlektionen am Game.\\
	Zwischenstand und Abgabe jeweils am Ende Ihrer Doppellektion.
\end{frame}

\begin{frame}{Zwischenstand}
	\begin{itemize}
		\item Sie zeigen eine spielbare Version und geben sie ab.
		\item Sie sehen, wo Ihre Gruppe steht, und ich sehe, wo es Hilfe braucht.
		\item Die Zwischenversion dient als Sicherung, falls bei der Abgabe etwas fehlt.
	\end{itemize}
\end{frame}

\begin{frame}{Bewertung}
	\begin{enumerate}
		\item \textbf{Produkt (Gruppe):} Ich bewerte das abgegebene Game nach Funktion, Code und Umfang.
		\item \textbf{Einzelgespräch (individuell):} Sie erklären eine Stelle im Code, die ich auswähle, oder ändern sie direkt.
	\end{enumerate}
	\vspace{1em}
	Wer den eigenen Code versteht, ist im Gespräch gut vorbereitet.
\end{frame}

\begin{frame}{Regeln}
	\begin{itemize}
		\item Der Abgabetermin ist fest.
		\item Die Einzelgespräche finden in Ihrer Halbklasse statt. Alle Mitglieder einer Gruppe müssen in derselben Halbklasse sein.
		\item Wer bei einem Einzelgespräch fehlt, holt es am zweiten Gesprächstag nach.
		\item Die Noten müssen vor dem Notenkonvent stehen. Später ist nichts mehr möglich.
	\end{itemize}
\end{frame}
